% =====================================================================
%  1.4 Responsabilidades das instituições
% =====================================================================
\section{Responsabilidades das instituições}
\label{sec:responsabilidades}

O IAT é o executor da ação: planeja, contrata, fiscaliza e valida os produtos.
O IDR-Paraná e a Sanepar são coexecutores. O IDR-Paraná atua na mobilização e
na frente dos mananciais, e a Sanepar executa o Programa Sanepar Rural com
recursos próprios. A SECID valida a implantação dos modelos de gestão. As
prefeituras participam da governança e da articulação local, e as comunidades,
da gestão e da operação dos sistemas. Empresas contratadas executam os
serviços e as obras, e a UGP consolida o monitoramento e o reporte ao Banco
Mundial \cite{mop2026}. O \autoref{qua:roteiro} apresenta o roteiro de cada
ação, e o \autoref{qua:matriz-responsabilidades} resume quem executa, apoia e
valida.

\begin{landscape}
\begin{quadro}[p]
\caption{Roteiro por ação: entradas, responsáveis e órgãos envolvidos}
\label{qua:roteiro}
\scriptsize
\renewcommand{\arraystretch}{1.25}
\begin{tabularx}{\linewidth}{|P{3.6cm}|L|P{3.2cm}|P{3.9cm}|L|}
\hline
\cab{Ação} & \cab{Entradas} & \cab{Responsável pela execução} &
\cab{Órgão de apoio ou validação} & \cab{Saída} \\ \hline
Planejamento integrado & Critérios do programa; diagnóstico estadual do saneamento rural &
  IAT (equipe interna) & IDR-Paraná & Planejamento das ações \\ \hline
Seleção dos municípios e comunidades & Critérios de seleção; Censo 2022; IPDM &
  IAT (equipe interna) & IDR-Paraná & Mapa dos municípios \\ \hline
Adesão e Termo de Cooperação & Mapa dos municípios; minuta do termo &
  IAT & Prefeituras & Lista de municípios; termos assinados \\ \hline
Definição e implantação do modelo de gestão & Arranjos de gestão em análise; Termo de Cooperação &
  IAT & SECID (validação); IDR-Paraná; Sanepar; prefeituras; comunidades &
  Modelo de gestão; associações regularizadas \\ \hline
Trabalho técnico-social & Termo de Cooperação; TR &
  Empresa contratada & IAT e IDR-Paraná (validação); prefeituras &
  Cadastro, diagnóstico e termos de adesão \\ \hline
Captação: poço e outorga & Diagnóstico; estudo hidrogeológico; materiais de perfuração &
  Empresa contratada & IAT (fiscalização e outorga); prefeituras & Poço perfurado e outorgado \\ \hline
Projetos e licenciamento & Vazão e qualidade da água; anteprojeto &
  Empresa projetista & IAT (aprovação e licenças) & Projetos aprovados; licenças \\ \hline
Obras de abastecimento & Projetos; materiais; licenças &
  Executor da obra & IAT (fiscalização); empresa supervisora; prefeituras &
  Sistema implantado; \textit{as built} \\ \hline
Supervisão e gerenciamento das obras & Contrato de supervisão; cronograma das obras &
  Empresa supervisora & IAT & Relatórios de supervisão \\ \hline
Esgotamento nas comunidades & Diagnóstico domiciliar; estudo de alternativas &
  Empresa contratada & IAT (validação); prefeituras; proprietários & Sistemas instalados e pagos \\ \hline
Esgotamento na área do IDR-Paraná & Planos Integrados de Propriedade (PIP) &
  Empresa contratada & IAT e IDR-Paraná (validação); proprietários & Sistemas instalados e pagos \\ \hline
Programa Sanepar Rural & Contrato de concessão; Termo de Compromisso &
  Sanepar & Prefeituras (obras); comunidades & Sistemas implantados \\ \hline
Entrega e operação & Sistema concluído; associação regularizada &
  IAT & Entidade gestora; prefeituras & Aceite da entidade gestora; TRD \\ \hline
Monitoramento e reporte & Relatórios; dados do SIGARH e GeoPR &
  IAT e Sanepar & UGP; Banco Mundial & Relatórios semestrais e anuais \\ \hline
\end{tabularx}
\fonte{Elaborado pelo IAT (2026) com base em \citeonline{mop2026} e no POA 2026--2027.}
\end{quadro}
\end{landscape}

\begin{quadro}[htbp]
\caption{Matriz de responsabilidades por ação}
\label{qua:matriz-responsabilidades}
\scriptsize
\renewcommand{\arraystretch}{1.3}
\newcommand{\rE}{\cellcolor{azulpsh}\color{white}\textbf{E}}
\newcommand{\rA}{\cellcolor{azulclaro}A}
\newcommand{\rV}{\cellcolor{orange!45}\textbf{V}}
\begin{tabularx}{\textwidth}{|L|c|c|c|c|c|c|c|c|}
\hline
\cab{Ação} & \cab{IAT} & \cab{IDR} & \cab{Sanepar} & \cab{SECID} &
\cab{Pref.} & \cab{Com.} & \cab{Empr.} & \cab{UGP} \\ \hline
Planejamento e seleção          & \rE & \rA &    &    &    &    &    &    \\ \hline
Adesão e Termo de Cooperação    & \rE &    &    &    & \rE &    &    &    \\ \hline
Modelo de gestão                & \rE & \rA & \rA & \rV & \rA & \rA &    &    \\ \hline
Trabalho técnico-social         & \rV & \rV &    &    & \rA & \rA & \rE &    \\ \hline
Captação, projetos e licenças   & \rV &    &    &    & \rA &    & \rE &    \\ \hline
Obras e supervisão              & \rV &    &    &    & \rA &    & \rE &    \\ \hline
Esgotamento -- comunidades      & \rV &    &    &    & \rA & \rA & \rE &    \\ \hline
Esgotamento -- área do IDR      & \rV & \rV &    &    &    & \rA & \rE &    \\ \hline
Programa Sanepar Rural          &    &    & \rE &    & \rE & \rA &    &    \\ \hline
Entrega e operação              & \rE &    &    &    & \rA & \rE &    &    \\ \hline
Monitoramento e reporte         & \rE &    & \rE &    &    &    &    & \rV \\ \hline
\end{tabularx}
\fonte{Elaborado pelo IAT (2026) com base em \citeonline{mop2026}.}
\nota{E = executa; A = apoia; V = valida ou fiscaliza. Pref.~=~prefeituras;
Com.~=~comunidades; Empr.~=~empresas contratadas.}
\end{quadro}

\pendencia{O MOP atribui ao IDR-Paraná corresponsabilidade pelo abastecimento
(Quadro 30: ``IAT/IDR''), mas não descreve o papel do instituto nessa frente
além da mobilização social. Detalhar.}
